\documentclass[a4paper,11pt]{article}
\usepackage[utf8]{inputenc}
\usepackage[T1]{fontenc}
\usepackage[french]{babel}
\usepackage{amsmath, amssymb, amsthm}
\usepackage{graphicx}
\usepackage{tcolorbox}
\usepackage{tikz}
\usetikzlibrary{shapes,arrows,positioning}
\usepackage{listings}
\usepackage{xcolor}
\usepackage{geometry}
\usepackage{hyperref}

\geometry{hmargin=2.5cm,vmargin=2.5cm}

% Configuration des boîtes
\tcbuselibrary{skins,breakable}

\newtcolorbox{conceptbox}[1][]{
  colback=blue!5!white,
  colframe=blue!75!black,
  title={\textbf{Concept Clé}},
  fonttitle=\bfseries,
  enhanced,
  drop shadow,
  breakable,
  #1
}

\newtcolorbox{mathbox}[1][]{
  colback=green!5!white,
  colframe=green!60!black,
  title={\textbf{Fondement Mathématique}},
  fonttitle=\bfseries,
  enhanced,
  drop shadow,
  breakable,
  #1
}

\newtcolorbox{codebox}[1][]{
  colback=gray!5!white,
  colframe=gray!75!black,
  title={\textbf{Implémentation Python}},
  fonttitle=\bfseries,
  breakable,
  #1
}

\newtcolorbox{warningbox}[1][]{
  colback=red!5!white,
  colframe=red!75!black,
  title={\textbf{Attention}},
  fonttitle=\bfseries,
  enhanced,
  #1
}

\newtcolorbox{examplebox}[1][]{
  colback=orange!5!white,
  colframe=orange!75!black,
  title={\textbf{Exemple}},
  fonttitle=\bfseries,
  enhanced,
  breakable,
  #1
}

% Configuration du code
\definecolor{codegreen}{rgb}{0,0.6,0}
\definecolor{codegray}{rgb}{0.5,0.5,0.5}
\definecolor{codepurple}{rgb}{0.58,0,0.82}
\definecolor{backcolour}{rgb}{0.95,0.95,0.92}

\lstdefinestyle{pythonstyle}{
    language=Python,
    backgroundcolor=\color{backcolour},   
    commentstyle=\color{codegreen},
    keywordstyle=\color{magenta},
    numberstyle=\tiny\color{codegray},
    stringstyle=\color{codepurple},
    basicstyle=\ttfamily\footnotesize,
    breakatwhitespace=false,         
    breaklines=true,                 
    captionpos=b,                    
    keepspaces=true,                 
    numbers=left,                    
    numbersep=5pt,                  
    showspaces=false,                
    showstringspaces=false,
    showtabs=false,                  
    tabsize=2
}
\lstset{style=pythonstyle}

\title{\textbf{t-SNE} \\ \large t-distributed Stochastic Neighbor Embedding \\ Visualisation de Données Haute Dimension}
\author{Cours de Machine Learning Non Supervisé}
\date{\today}

\begin{document}

\maketitle
\tableofcontents
\newpage

\section{Introduction}

Le \textbf{t-SNE} (t-distributed Stochastic Neighbor Embedding) est une technique de réduction de dimensionnalité développée par Laurens van der Maaten et Geoffrey Hinton en 2008.

\begin{conceptbox}
\textbf{Objectif Principal}

Contrairement à la PCA qui préserve les distances globales, t-SNE cherche à :
\begin{itemize}
    \item Préserver les \textbf{relations de voisinage local}
    \item Projeter des données haute dimension en 2D ou 3D pour \textbf{visualisation}
    \item Révéler la \textbf{structure en clusters} cachée dans les données
\end{itemize}

t-SNE est principalement un outil de \textbf{visualisation}, pas de réduction de dimension pour le machine learning.
\end{conceptbox}

\section{Pourquoi t-SNE après l'Espace Latent ?}

\subsection{Le Problème de Visualisation}

Après avoir entraîné un autoencodeur, vous obtenez un espace latent de dimension $d$ (par exemple, $d=2$ ou $d=32$).

\begin{examplebox}
\textbf{Cas Pratique : Votre Autoencodeur}

\begin{itemize}
    \item \textbf{Entrée} : 5 features (température, vitesse, couple, usure)
    \item \textbf{Espace latent} : 2 dimensions
    \item \textbf{Problème résolu ?} Oui, vous pouvez déjà visualiser en 2D !
\end{itemize}

\textbf{Mais si votre espace latent avait 32 dimensions ?}
\begin{itemize}
    \item Impossible de visualiser directement
    \item t-SNE permet de passer de 32D à 2D tout en préservant les clusters
\end{itemize}
\end{examplebox}

\subsection{Quand Utiliser t-SNE ?}

\begin{center}
\begin{tabular}{|l|c|c|}
\hline
\textbf{Situation} & \textbf{t-SNE Utile ?} & \textbf{Raison} \\
\hline
Espace latent 2D & Non & Déjà visualisable \\
Espace latent 3D & Optionnel & Peut améliorer la séparation \\
Espace latent $>$ 3D & Oui & Indispensable pour visualiser \\
\hline
\end{tabular}
\end{center}

\section{Principe de Fonctionnement}

\subsection{Intuition}

t-SNE fonctionne en deux étapes :
\begin{enumerate}
    \item \textbf{Espace original} : Calculer les probabilités de voisinage entre tous les points
    \item \textbf{Espace réduit} : Placer les points en 2D de façon à reproduire ces probabilités
\end{enumerate}

\subsection{Algorithme Simplifié}

\begin{mathbox}
\textbf{Étape 1 : Similarités dans l'Espace Original}

Pour chaque paire de points $(x_i, x_j)$ en haute dimension, on calcule :
$$ p_{j|i} = \frac{\exp(-||x_i - x_j||^2 / 2\sigma_i^2)}{\sum_{k \neq i} \exp(-||x_i - x_k||^2 / 2\sigma_i^2)} $$

Puis on symétrise :
$$ p_{ij} = \frac{p_{j|i} + p_{i|j}}{2n} $$

\textbf{Interprétation :} $p_{ij}$ mesure la probabilité que $x_i$ choisisse $x_j$ comme voisin.
\end{mathbox}

\begin{mathbox}
\textbf{Étape 2 : Similarités dans l'Espace Réduit (2D)}

Pour les points projetés $(y_i, y_j)$ en 2D, on utilise une distribution de Student (t) :
$$ q_{ij} = \frac{(1 + ||y_i - y_j||^2)^{-1}}{\sum_{k \neq l} (1 + ||y_k - y_l||^2)^{-1}} $$

\textbf{Pourquoi la distribution t ?}
\begin{itemize}
    \item Elle a des queues plus lourdes qu'une gaussienne
    \item Permet de mieux séparer les clusters en 2D
    \item Évite le "crowding problem" (entassement des points)
\end{itemize}
\end{mathbox}

\begin{mathbox}
\textbf{Étape 3 : Optimisation}

On minimise la divergence de Kullback-Leibler entre $P$ et $Q$ :
$$ \text{KL}(P||Q) = \sum_{i \neq j} p_{ij} \log \frac{p_{ij}}{q_{ij}} $$

Cette optimisation se fait par descente de gradient.
\end{mathbox}

\section{Hyperparamètres Clés}

\subsection{Perplexité (Perplexity)}

C'est le paramètre le plus important.

\begin{itemize}
    \item \textbf{Définition} : Mesure du nombre effectif de voisins considérés pour chaque point
    \item \textbf{Valeur typique} : 5 à 50 (défaut : 30)
    \item \textbf{Effet} :
    \begin{itemize}
        \item Perplexité faible (5-10) : Focus sur la structure locale (petits clusters)
        \item Perplexité élevée (30-50) : Capture la structure globale
    \end{itemize}
\end{itemize}

\begin{warningbox}
\textbf{Règle Empirique}

La perplexité doit être inférieure au nombre de points :
$$ 5 \leq \text{perplexity} < \frac{n}{3} $$

Pour $n=10000$ points, tester : 30, 50, 100.
\end{warningbox}

\subsection{Nombre d'Itérations}

\begin{itemize}
    \item \textbf{Défaut} : 1000 itérations
    \item \textbf{Recommandation} : Au moins 1000, souvent 2000-5000 pour convergence
    \item \textbf{Indicateur} : Surveiller la valeur de KL divergence (doit décroître)
\end{itemize}

\subsection{Learning Rate}

\begin{itemize}
    \item \textbf{Défaut} : 200
    \item \textbf{Plage} : 10 à 1000
    \item Si les points s'agglomèrent en boule : augmenter
    \item Si les points explosent : diminuer
\end{itemize}

\section{Avantages et Limites}

\subsection{Avantages}

\begin{itemize}
    \item Excellente préservation de la structure locale
    \item Révèle des clusters non linéaires
    \item Visualisations très intuitives
\end{itemize}

\subsection{Limites Critiques}

\begin{warningbox}
\textbf{Limitations Importantes}

\begin{enumerate}
    \item \textbf{Non déterministe} : Deux exécutions donnent des résultats différents
    \item \textbf{Distances non préservées} : Les distances en 2D n'ont PAS de sens absolu
    \item \textbf{Lent} : Complexité $O(n^2)$, difficile au-delà de 10 000 points
    \item \textbf{Pas de prédiction} : Ne peut pas projeter de nouvelles données
\end{enumerate}
\end{warningbox}

\section{Exemple de Code Python}

\begin{codebox}
\begin{lstlisting}
from sklearn.manifold import TSNE
import matplotlib.pyplot as plt
import numpy as np

# Supposons que Z est votre espace latent (n_samples, n_features)
# Par exemple, Z de dimension (10000, 32)

# 1. CONFIGURATION
tsne = TSNE(
    n_components=2,      # Projection en 2D
    perplexity=30,       # Nombre de voisins effectifs
    n_iter=1000,         # Nombre d'itérations
    random_state=42,     # Pour reproductibilité
    verbose=1            # Affiche la progression
)

# 2. PROJECTION
print("Calcul du t-SNE en cours...")
Z_tsne = tsne.fit_transform(Z)

# 3. VISUALISATION
plt.figure(figsize=(10, 8))

# Colorer par les vraies étiquettes (si disponibles)
scatter = plt.scatter(Z_tsne[:, 0], Z_tsne[:, 1], 
                     c=y, cmap='viridis', 
                     alpha=0.6, s=10)
plt.colorbar(scatter, label='Machine Failure')
plt.title('Visualisation t-SNE de l\'Espace Latent')
plt.xlabel('t-SNE Dimension 1')
plt.ylabel('t-SNE Dimension 2')
plt.grid(True, alpha=0.3)
plt.show()

# 4. ANALYSE PAR CLUSTER (si clustering fait)
# Exemple avec les clusters K-Means
plt.figure(figsize=(10, 8))
for cluster_id in np.unique(clusters):
    mask = clusters == cluster_id
    plt.scatter(Z_tsne[mask, 0], Z_tsne[mask, 1], 
               label=f'Cluster {cluster_id}', alpha=0.6, s=10)
plt.legend()
plt.title('t-SNE coloré par Clusters K-Means')
plt.show()
\end{lstlisting}
\end{codebox}

\section{Cas d'Usage dans Votre Projet}

\subsection{Votre Situation Actuelle}

\begin{itemize}
    \item Espace latent : \textbf{2 dimensions} (déjà visualisable)
    \item Clusters K-Means : 3 clusters identifiés
    \item Clusters DBSCAN : 2 clusters identifiés
\end{itemize}

\textbf{Conclusion :} t-SNE n'est \textbf{pas nécessaire} pour votre cas actuel.

\subsection{Quand Serait-il Utile ?}

Si vous aviez un espace latent de dimension supérieure (ex: 32D), t-SNE permettrait de :
\begin{enumerate}
    \item Visualiser la séparation entre Normal/Dégradé/Panne
    \item Vérifier si les clusters trouvés par K-Means sont bien séparés visuellement
    \item Identifier des sous-structures non détectées
\end{enumerate}

\section{Conclusion}

t-SNE est un outil puissant de \textbf{visualisation exploratoire} pour comprendre la structure de données haute dimension. Cependant :

\begin{itemize}
    \item Ne l'utilisez \textbf{pas} pour du machine learning (pas de prédiction)
    \item Ne l'utilisez \textbf{pas} si votre espace est déjà 2D ou 3D
    \item Utilisez-le pour \textbf{explorer} et \textbf{communiquer} vos résultats
\end{itemize}

Pour votre projet avec un espace latent 2D, vous pouvez directement visualiser $Z$ sans t-SNE.

\end{document}
